% ============================================================
%  main.tex — 具身智能 SCI 论文骨架(期刊无关)
%  用法:cp -R references/latex-skeleton/* reports/drafts/paper-<id>/
%        再把 \documentclass 换成目标期刊官方模板(见 references/journal-templates.md)
%  纪律:每个数字回指 ⑥ 审计、每个主张回指来源;未审计数字用 \audittodo{} 占位,禁写死。
% ============================================================

% --- RA-L / T-RO 用 IEEEtran;IJRR 用 SAGE;按期刊替换下面这行 ---
\documentclass[journal]{IEEEtran}

\usepackage{amsmath,amssymb,amsfonts}
\usepackage{graphicx}
\usepackage{booktabs}          % 三线表
\usepackage{siunitx}           % 单位与数字对齐
\usepackage{algorithm}
\usepackage{algorithmic}
\usepackage{xcolor}
\usepackage{hyperref}
\usepackage{subcaption}        % 定性图/失败案例并排
\usepackage[numbers,sort&compress]{natbib} % IEEEtran 数字引用;按期刊调整

% --- 诚实闸宏:未过 ⑥ 审计的数字一律占位,不得写死 ---
\newcommand{\audittodo}[1]{\textcolor{red}{[AUDIT-TODO: #1]}}
% dry-run:⑥ 对 exp-VP01 签发 ❌(数据 PLACEHOLDER)→ 全文数字一律 \audittodo{},禁写死。
% 真实运行示例格式:成功率 \num{XX.X}\% \pm \num{Y.Y}  % <- tests/exp-<id>, ⑥ 签发✅ 日期, 配置口径

\graphicspath{{figures/}}

\begin{document}

% 🔴 署名/单位由人确定,AI 不擅自填写 —— 保留占位,提交前须委托人明示
\title{Paper Title Here\\ \vspace{2pt}
\large (working title — 由委托人/作者定稿)}
\author{%
  \IEEEauthorblockN{Author One\IEEEauthorrefmark{1}, Author Two\IEEEauthorrefmark{2}}
  \IEEEauthorblockA{\IEEEauthorrefmark{1}Affiliation One — \textcolor{red}{[署名待委托人确认]}}
  \IEEEauthorblockA{\IEEEauthorrefmark{2}Affiliation Two — \textcolor{red}{[署名待委托人确认]}}
}
\maketitle

\begin{abstract}
% ============================================================
%  abstract.tex — 摘要(按期刊词数上限:RA-L ~150 词 / T-RO 可稍长)
%  纪律:数字与正文严格一致(最高频事故源);一句话讲清 问题-方法-结果-意义。
%  未审计数字用 \audittodo{} 占位,禁写死。
% ============================================================

% [问题] 一句话说清要解决的具身智能问题与其难点(不铺陈大背景)。
Long-horizon dexterous manipulation remains difficult because \ldots

% [方法] 一句话概括提出的方法及其核心创新(可证伪,不夸大)。
We propose \textbf{MethodName}, which \ldots

% [结果] 关键定量结果 —— 每个数字必须与正文/表格逐字一致,且带配置口径。
% 措辞纪律(⑥§2 + ⑨F-越级修正):配对 n=5 检验力不足(p 下限 0.0625)→ 禁用比较级/显著性声称。
On \textsc{LIBERO-Pro} frozen subset (single-arm, $N{=}300$), the three arms
differ descriptively: \audittodo{exp-VP01 三臂成功率, ⑥ 签❌占位}; with $n{=}5$
paired seeds the margin is \emph{not} statistically significant (permutation
$p$ floor $=0.0625$), so we report it as descriptive and defer significance to
additional seeds. % 承 ⑥ 降级:不声称显著、不用比较级(检验力不足)

% [意义] 一句话说明为何重要 / 对领域的可操作启示(不超出证据)。
These results suggest \ldots Code and data are available at \url{<repo, 提交前由人确认>}.

\end{abstract}

\begin{IEEEkeywords}
manipulation, dexterous hands, robot learning, \ldots % 按期刊关键词数
\end{IEEEkeywords}

% ============================================================
%  intro.tex — 引言:motivation → gap → 贡献点(编号,可证伪)
%  纪律:主张回指来源(见 traceability-worksheet.md 主张台账);比较级须有支撑。
% ============================================================

\section{Introduction}
\label{sec:intro}

% [动机] 从领域真实痛点切入,不空泛。每条事实性主张 \cite{}。
Dexterous manipulation is central to embodied intelligence because \ldots~\cite{<key>}.
However, current approaches struggle with \ldots~\cite{<key>}.

% [gap] 明确指出现有工作的具体缺口(与 ① GapScout 的差异化陈述一致)。
Prior work~\cite{<key>,<key>} addresses \ldots, yet \textbf{remains limited in}
\ldots. This gap matters because \ldots

% [本文思路] 一句话点出核心 insight / 方法直觉。
In this paper, we argue that \ldots and propose \textbf{MethodName}.

% [贡献点] 编号、可证伪、每条对应后文实验。禁「首个/最优」等无支撑最高级。
Our contributions are:
\begin{itemize}
  \item \textbf{C1:} A method that \ldots (validated in Sec.~\ref{sec:exp}, Table~\ref{tab:main}).
  \item \textbf{C2:} An ablation showing \ldots (Sec.~\ref{sec:ablation}).
  \item \textbf{C3:} Real-robot results demonstrating \ldots (Sec.~\ref{sec:real}, 采集条件见附录).
\end{itemize}

% [路线图] 可选:简述后文结构。

% ============================================================
%  related.tex — 相关工作:批判性归类(引 ② taxonomy),不罗列堆砌
%  纪律:每条引用来自打开过原文的语料清单;说清本工作与各路线的差异。
% ============================================================

\section{Related Work}
\label{sec:related}

% 按 ② SurveyWeaver 的 method taxonomy 分小节,每节末点出「与本工作的差异」。

\subsection{Route A: <名称, 如 Imitation Learning for Manipulation>}
Representative works~\cite{<key>,<key>} \ldots. Their strength is \ldots,
but they \textbf{differ from ours in} \ldots.

\subsection{Route B: <名称, 如 Tactile / Multi-modal Sensing>}
\ldots~\cite{<key>}. In contrast, we \ldots

\subsection{Route C: <名称, 如 Language-conditioned / VLA Policies>}
\ldots~\cite{<key>}. Unlike these, our approach \ldots

% [定位小结] 一段话把本工作放进版图,明确差异化(回应 ① 的 gap 陈述)。
\paragraph{Positioning.} Compared with the above, MethodName is distinct in that
\ldots, which directly addresses the gap identified in Sec.~\ref{sec:intro}.

% ============================================================
%  method.tex — 方法:可复现细节充足,读者能照做
%  纪律:公式符号一致;架构/损失/超参齐全;区分「方法本身」与「实现选择」。
% ============================================================

\section{Method}
\label{sec:method}

% [问题形式化] 记号定义(POMDP / 策略 π(a|o) 等),符号全篇一致。
\subsection{Problem Formulation}
We consider \ldots. Let $o_t$ denote \ldots and $a_t$ denote \ldots.
The goal is to learn a policy $\pi_\theta(a_t \mid o_t, \ell)$ conditioned on
language instruction $\ell$ \ldots

% [方法总览] 配架构图(见 ⑦ figure-builder, diagram 类)。
\subsection{Overview}
Figure~\ref{fig:pipeline} illustrates MethodName, which consists of \ldots
\begin{figure}[t]
  \centering
  \includegraphics[width=\linewidth]{pipeline}
  \caption{Overview of MethodName. \ldots (caption 须自解释)}
  \label{fig:pipeline}
\end{figure}

% [核心组件] 逐个讲清提出的模块;每个组件对应后文一个消融。
\subsection{Component 1: <名称>}
\ldots
\begin{equation}
  \mathcal{L} = \ldots
  \label{eq:loss}
\end{equation}

\subsection{Component 2: <名称>}
\ldots

% [训练/推理] 算法伪代码(algorithm 环境);关键超参见附录,不散落正文。
\subsection{Training and Inference}
\begin{algorithm}[t]
\caption{MethodName Training}
\begin{algorithmic}[1]
\STATE \ldots
\end{algorithmic}
\end{algorithm}

% [实现选择 vs 方法本质] 明确哪些是可替换的工程选择,避免过度声称普适性。
\paragraph{Design choices.} We emphasize that \ldots is an implementation choice,
whereas \ldots is central to the method.

% ============================================================
%  experiments.tex — R3-NeuTraL(dry-run 验证;数字全 \audittodo 占位)
%  纪律:⑥ ResultsAuditor 对 exp-VP01 签发 ❌(数据 PLACEHOLDER)→ 正文数字
%        一律 \audittodo{},禁写死(呼应 paper-writer §1 前置门 + ⑥§9 放行去向)。
% ============================================================

\section{Experiments}
\label{sec:exp}

\subsection{Experimental Setup}
\paragraph{Benchmark.} We evaluate on \textsc{LIBERO-Pro}~\cite{liberopro2025}
(official 1{,}648 perturbation variants), using a \emph{stratified frozen subset}
of $N{=}300$ variants (sampling seed 20261009, rules public), \emph{single-arm}
(Franka Panda) configuration. % 配置口径:单臂,与 ⑥§4 一致,防跨配置迁移
\paragraph{Three arms (equal, not baseline-subordinate).}
(A) end-to-end VLA ($\pi_{0.5}$ frozen, open-weights);
(B) agentic harness (version-locked frontier VLM);
(C) code-as-policy (same VLM planner + public primitive library).
\paragraph{Protocol.} 50 episodes/task (official protocol); unified budget
($\leq$8 LLM calls/episode); metrics: success rate, recovery rate, first-failure
step, recovery latency, task-wipeout rate, cost (LLM calls/tokens/wallclock).
\paragraph{Seeds.} $N{=}\,$\audittodo{exp-VP01 seed 数} seeds; mean $\pm$ std.

\subsection{Main Results}
\begin{table}[t]
\centering
\caption{Three-arm neutral comparison on \textsc{LIBERO-Pro} frozen subset
(single-arm, $N{=}300$). \textbf{All numbers PLACEHOLDER (dry-run); audit signed
\texttt{NOT-CLEARED}.} Best in bold.}
\label{tab:main}
\begin{tabular}{lccc}
\toprule
Arm & Success (\%)$\uparrow$ & Recovery (\%) & Cost (LLM calls) \\
\midrule
(A) end-to-end VLA (frozen)   & \audittodo{exp-VP01/R3_vla_frozen} & \audittodo{} & \audittodo{} \\
(B) agentic harness           & \audittodo{exp-VP01/R3_harness}    & \audittodo{} & \audittodo{} \\
(C) code-as-policy            & \audittodo{exp-VP01/R3_cap}        & \audittodo{} & \audittodo{} \\
\bottomrule
\end{tabular}
\end{table}
% 措辞纪律(⑥§2):配对 n=5 检验力不足 → 禁写「B 显著优于 A」,即使占位差 +26.8pp。
Arm (B) shows a descriptive margin over (A) of \audittodo{delta}; however, with
$n{=}5$ paired seeds the permutation test cannot reach $p<0.05$ (min achievable
$p=0.0625$), so we \emph{do not} claim statistical significance and defer to
additional seeds (see Sec.~\ref{sec:limit}).

\subsection{Baseline-Alignment Check}
\label{sec:baseline}
The same frozen $\pi_{0.5}$ yields \audittodo{exp-VP01/R3_vla_frozen} under our
caliber vs.\ \audittodo{exp-VP01/baseline_align_pi05} under the alignment probe,
against literature values 12.8\% / 17.5\% / $\leq$13.3\% (Pigey / DynaHarness /
MotorMind, verified in Sec.~\ref{sec:related}). If $\geq$2 cannot be
simultaneously approached, ``even the baseline is not comparable'' is confirmed
--- the hardest evidence for the caliber-confusion thesis (Fig.~\ref{fig:spread}).

\subsection{Caliber Confusion (real data)}
Across 15 recent \textsc{LIBERO-Pro} reports the absolute success rate spans
12.8\%--95.0\% with $\geq$9 mutually incompatible denominators
(Fig.~\ref{fig:spread}); \emph{no two reports share a denominator}. These 15
values are externally verified (arXiv/OpenReview, 2026-10-09), not placeholder.

\subsection{Recovery and Cost Axes}
\label{sec:recovery}
Arms (B)/(C) show \audittodo{recovery vs A} recovery rate but
\audittodo{latency} higher recovery latency; on object/position static
perturbations the recovery gap \audittodo{vanishes?}. % 占位,待真实

\subsection{Qualitative and Failure Cases}
Representative rollouts and failure modes are shown in Fig.~\ref{fig:qual}
(placeholder frames; real capture pending stage-4 execution).

\subsection{Limitations}
\label{sec:limit}
(i) Cross-arm ``same backbone'' is physically infeasible (open-weight VLA vs
closed frontier VLM); ``aligned caliber'' is redefined as same-generation
strongest + public backbone identity + reported cost. (ii) $N{=}300$ frozen
subset $\neq$ official full 1{,}648; conclusions are scoped to the subset.
(iii) $n{=}5$ paired seeds are underpowered for significance claims.
(iv) Null-result risk: if budget effect $\gg$ route effect, the finding becomes
``route matters little, budget matters'' --- still falsifiable and publishable.

% ============================================================
%  conclusion.tex — 结论:不超出证据的总结 + 诚实的局限与未来工作
%  纪律:局限如实写(呼应 ⑤/④ 的失败案例与残差不确定性);不夸大普适性。
% ============================================================

\section{Conclusion}
\label{sec:conclusion}

% [总结] 一段话回收贡献点,措辞与结果强度匹配(不用超出证据的最高级)。
We presented \textbf{MethodName}, which \ldots. Across \textsc{Benchmark} and
real-robot experiments, it \ldots (Sec.~\ref{sec:exp}).

% [局限] 诚实条款落地:明说方法边界、未覆盖场景、已知失败模式。
\paragraph{Limitations.} MethodName currently \ldots; it has not been validated on
\ldots, and performance degrades when \ldots. These are honest boundaries of the
present study, not claims of generality.

% [未来工作] 可操作的下一步,不空泛。
\paragraph{Future work.} We plan to \ldots, and to extend to \ldots.


% --- 附录:复现材料(见 reproducibility.md) ---
\appendix
% ============================================================
%  appendix.tex — 附录:复现材料 + 额外结果 + 真机采集条件
%  纪律:实现细节(超参/版本)集中此处;真机采集条件须与 ⑤ 记录一致(🔴 人签名)。
% ============================================================

\section{Implementation Details}
\label{app:impl}
% 训练/评测的全部超参、硬件、仿真器与 benchmark 版本(呼应 ④ 冻结目录)。
\begin{itemize}
  \item Simulator: \ldots (version \ldots); GPU: \ldots
  \item Optimizer / lr / batch / epochs: \ldots
  \item Random seeds: \ldots; number of seeds $N=$ \audittodo{}
  \item Code commit: \texttt{<hash>}; config: \texttt{tests/exp-<id>/config.yaml}
\end{itemize}

\section{Additional Results}
\label{app:more}
% 完整结果表(含正文省略的、方差大的、不利的数字 —— 不 cherry-pick)。
\ldots

\section{Real-Robot Collection Details}
\label{app:real}
% 与 ⑤ realexec-checklist 的记录一致:硬件型号/固件、标定、遥操作方式、
% episode 数与成功判据、操作者签名确认的采集条件。🔴 由人回填并签名。
\begin{itemize}
  \item Hardware: arm \ldots (firmware \ldots), hand \ldots
  \item Calibration: hand-eye error \ldots (date \ldots)
  \item Teleoperation: \ldots; operator sign-off: \textcolor{red}{[待人签名]}
  \item Episodes: collected \ldots / success \ldots / failure \ldots
  \item Deviations from protocol: \ldots (如实写)
\end{itemize}

% --- 复现声明(见 reproducibility.md,按期刊要求放置) ---
\section*{Reproducibility Statement}
% ============================================================
%  reproducibility.tex — 复现声明(被 appendix.tex \input)
%  纪律:代码/数据/超参可获取性如实声明;开源链接提交前由人确认(🔴 对外发布)。
% ============================================================

% 若期刊用独立 reproducibility checklist(如 NeurIPS/RA-L),按其字段逐项填。
\textbf{Code.} Our training and evaluation code, including all configs used to
produce the reported numbers, is available at \url{<repo>} — \textcolor{red}{[开源链接与是否公开由委托人确认]}。

\textbf{Data.} Teleoperated / simulation data used in this work are
\ldots (available at \url{<link>} / not released due to \ldots). % 如实,不空承诺

\textbf{Hyperparameters.} Full hyperparameters, environment versions, and random
seeds are listed in Appendix~\ref{app:impl}; each reported number maps to a
frozen config under \texttt{tests/exp-<id>/}.

\textbf{Compute.} Experiments were run on \ldots (approx. \ldots GPU-hours).

\textbf{Reproducibility note.} Numbers in this paper were produced by frozen
configs (commit \texttt{<hash>}) and audited for full-caliber consistency
(including unfavorable results) before citation.



% --- 参考文献:来自 refs.bib(构建规范见 references/refs-bib-guide.md,零虚构) ---
\bibliographystyle{IEEEtran}
\bibliography{refs}

\end{document}
